\textbf{Phần 1. Trắc nghiệm lựa chọn}
\setcounterpageref{ex}{0}

\begin{ex} %[2H6V32]
Một ngày có $24$ giờ. Hỏi $3$ ngày có tất cả bao nhiêu giờ?
\choice
{$48$ giờ}
{\True $72$ giờ}
{$60$ giờ}
{$96$ giờ}
\loigiai{
$24 + 24 + 24 = 72$ (giờ).
\textbf{Đáp án của câu hỏi là B}
}
\end{ex}

\begin{ex} %[2H6V32]
Đồng hồ chỉ $7$ giờ $15$ phút tối. Cách đọc nào sau đây theo hệ $24$ giờ là đúng?
\choice
{$17$ giờ $15$ phút}
{$18$ giờ $15$ phút}
{$20$ giờ $15$ phút}
{\True $19$ giờ $15$ phút}
\loigiai{
$7$ giờ tối $= 19$ giờ. Vậy là $19$ giờ $15$ phút.
\textbf{Đáp án của câu hỏi là D}
}
\end{ex}

\begin{ex} %[2H6V32]
Tháng nào sau đây có đúng $30$ ngày?
\choice
{\True Tháng 4}
{Tháng 3}
{Tháng 5}
{Tháng 7}
\loigiai{
Tháng 4 có đúng 30 ngày.
\textbf{Đáp án của câu hỏi là A}
}
\end{ex}

\begin{ex} %[2H6V32]
Nếu hôm nay là thứ Năm ngày $8$. Hỏi $10$ ngày nữa là ngày mấy?
\choice
{Ngày $17$}
{Ngày $19$}
{\True Ngày $18$}
{Ngày $16$}
\loigiai{
$8 + 10 = 18$.
\textbf{Đáp án của câu hỏi là C}
}
\end{ex}

\begin{ex} %[2H6V32]
Bác thợ mộc đóng một chiếc bàn mất $3$ giờ. Bác bắt đầu đóng lúc $8$ giờ sáng và nghỉ giải lao $30$ phút. Bác đóng xong chiếc bàn lúc mấy giờ?
\choice
{$11$ giờ}
{\True $11$ giờ $30$ phút trưa}
{$12$ giờ}
{$10$ giờ $30$ phút}
\loigiai{
Thời gian hoàn thành: $8$ giờ $+ 3$ giờ $+ 30$ phút $= 11$ giờ $30$ phút.
\textbf{Đáp án của câu hỏi là B}
}
\end{ex}

\begin{ex} %[2H6V32]
Kim giờ quay được $1$ vòng thì kim phút quay được bao nhiêu vòng?
\choice
{$6$ vòng}
{$24$ vòng}
{$60$ vòng}
{\True $12$ vòng}
\loigiai{
Kim giờ quay 1 vòng hết 12 giờ. Mỗi giờ kim phút quay 1 vòng nên quay được 12 vòng.
\textbf{Đáp án của câu hỏi là D}
}
\end{ex}

\begin{ex} %[2H6V32]
Một tháng có $31$ ngày, nếu ngày đầu tiên là thứ Ba thì tháng đó có bao nhiêu ngày thứ Ba?
\choice
{\True $5$ ngày}
{$4$ ngày}
{$6$ ngày}
{$3$ ngày}
\loigiai{
Các ngày thứ Ba là: 1, 8, 15, 22, 29 (có 5 ngày).
\textbf{Đáp án của câu hỏi là A}
}
\end{ex}

\begin{ex} %[2H6V32]
Tàu hỏa chạy từ ga $A$ lúc $6$ giờ sáng và đến ga $B$ lúc $14$ giờ cùng ngày. Tàu đã chạy hết bao nhiêu giờ?
\choice
{$7$ giờ}
{$9$ giờ}
{\True $8$ giờ}
{$10$ giờ}
\loigiai{
$14 - 6 = 8$ (giờ).
\textbf{Đáp án của câu hỏi là C}
}
\end{ex}

\begin{ex} %[2H6V32]
Có bao nhiêu ngày trong $4$ tuần lễ?
\choice
{$24$ ngày}
{\True $28$ ngày}
{$30$ ngày}
{$32$ ngày}
\loigiai{
$7 \times 4 = 28$ (ngày).
\textbf{Đáp án của câu hỏi là B}
}
\end{ex}

\begin{ex} %[2H6V32]
Đồng hồ chỉ $16$ giờ $30$ phút. Kim phút đang chỉ vào số nào trên mặt đồng hồ?
\choice
{Số $3$}
{Số $12$}
{Số $9$}
{\True Số $6$}
\loigiai{
$30$ phút ứng với kim phút chỉ vào số 6.
\textbf{Đáp án của câu hỏi là D}
}
\end{ex}

\begin{ex} %[2H6V32]
Nếu thứ Bảy là ngày $15$ thì Chủ nhật tuần sau là ngày mấy?
\choice
{\True Ngày $23$}
{Ngày $22$}
{Ngày $24$}
{Ngày $21$}
\loigiai{
Chủ nhật tuần này là ngày 16. Chủ nhật tuần sau là: $16 + 7 = 23$.
\textbf{Đáp án của câu hỏi là A}
}
\end{ex}

\begin{ex} %[2H6V32]
Mỗi ngày Nam dành $1$ giờ để đọc sách. Hỏi trong một tuần Nam dành bao nhiêu giờ đọc sách?
\choice
{$5$ giờ}
{$6$ giờ}
{\True $7$ giờ}
{$8$ giờ}
\loigiai{
Một tuần có 7 ngày: $1 \times 7 = 7$ (giờ).
\textbf{Đáp án của câu hỏi là C}
}
\end{ex}

\textbf{Phần 2. Trắc nghiệm đúng sai}
\setcounterpageref{ex}{0}

\begin{ex} %[2H6V32]
Xét tính đúng sai của các bài toán thời gian tổng hợp:
\choiceTF
{\True Một năm thông thường có 365 ngày.}
{Năm nào tháng 2 cũng có 29 ngày.}
{\True Khoảng thời gian từ 10 giờ sáng đến 2 giờ chiều là 4 giờ.}
{Một tuần lễ chỉ gồm các ngày học từ thứ Hai đến thứ Sáu.}
\loigiai{
\begin{itemize}
    \item Ý a) Đúng: Đúng theo lịch dương chuẩn.
    \item Ý b) Sai: Chỉ có năm nhuận tháng 2 mới có 29 ngày.
    \item Ý c) Đúng: 2 giờ chiều là 14 giờ: $14 - 10 = 4$ giờ. Đúng.
    \item Ý d) Sai: Một tuần lễ đầy đủ gồm 7 ngày từ thứ Hai đến Chủ nhật.
\end{itemize}
\begin{itemize}
    \item a) \textbf{Đúng}.
    \item b) \textbf{Sai}.
    \item c) \textbf{Đúng}.
    \item d) \textbf{Sai}.
\end{itemize}
}
\end{ex}

\begin{ex} %[2H6V32]
Kỳ nghỉ lễ Quốc khánh 2 tháng 9:
\choiceTF
{\True Tháng 9 có đúng 30 ngày.}
{\True Ngày 2 tháng 9 là ngày Quốc khánh của nước Việt Nam.}
{Nếu ngày 2 tháng 9 là thứ Hai thì ngày 9 tháng 9 là thứ Ba.}
{\True Nếu được nghỉ 4 ngày từ ngày 1 đến hết ngày 4 tháng 9 thì học sinh đi học lại vào ngày 5 tháng 9.}
\loigiai{
\begin{itemize}
    \item Ý a) Đúng: Đúng.
    \item Ý b) Đúng: Đúng.
    \item Ý c) Sai: Cộng 7 ngày nên ngày 9 tháng 9 vẫn là thứ Hai.
    \item Ý d) Đúng: Đúng.
\end{itemize}
\begin{itemize}
    \item a) \textbf{Đúng}.
    \item b) \textbf{Đúng}.
    \item c) \textbf{Sai}.
    \item d) \textbf{Đúng}.
\end{itemize}
}
\end{ex}

\textbf{Phần 3. Trả lời ngắn}
\setcounterpageref{ex}{0}

\begin{ex} %[2H6V32]
Một ngày có 24 giờ. Hỏi nửa ngày có bao nhiêu giờ?
\shortans{12}
\loigiai{
\textbf{Đáp số:} $12$\\
$24 : 2 = 12$ giờ.
}
\end{ex}

\begin{ex} %[2H6V32]
Nếu ngày 5 của một tháng là thứ Ba thì ngày thứ Ba của tuần tiếp theo là ngày mấy?
\shortans{12}
\loigiai{
\textbf{Đáp số:} $12$\\
$5 + 7 = 12$.
}
\end{ex}

\begin{ex} %[2H6V32]
Một tuần lễ có 7 ngày. Hỏi 5 tuần lễ có tất cả bao nhiêu ngày?
\shortans{35}
\loigiai{
\textbf{Đáp số:} $35$\\
$7 \times 5 = 35$ ngày.
}
\end{ex}

\begin{ex} %[2H6V32]
Bác thợ làm việc từ 7 giờ sáng đến 11 giờ trưa và từ 13 giờ đến 17 giờ chiều. Hỏi bác làm việc tất cả bao nhiêu giờ trong ngày?
\shortans{8}
\loigiai{
\textbf{Đáp số:} $8$\\
Sáng: $11 - 7 = 4$ giờ. Chiều: $17 - 13 = 4$ giờ. Tổng: $4 + 4 = 8$ giờ.
}
\end{ex}

\begin{ex} %[2H6V32]
Đồng hồ chỉ 21 giờ 30 phút, tức là 9 giờ bao nhiêu phút tối?
\shortans{30}
\loigiai{
\textbf{Đáp số:} $30$\\
30 phút.
}
\end{ex}

\begin{ex} %[2H6V32]
Tháng Mười có 31 ngày, tháng Mười Một có 30 ngày. Hỏi hai tháng đó có tất cả bao nhiêu ngày?
\shortans{61}
\loigiai{
\textbf{Đáp số:} $61$\\
$31 + 30 = 61$ ngày.
}
\end{ex}
